\section{Sujet n\textsuperscript{o}\,10}
\noindent\begin{minipage}{0.5\linewidth}\footnotesize Office du Baccalauréat du Cameroun\end{minipage}%
\hfill\begin{minipage}{0.45\linewidth}\footnotesize\raggedleft Examen : \textbf{Baccalauréat} · Série : \textbf{C/E}\\
Épreuve : \textbf{Mathématiques} · Durée : \textbf{4\,h} · Coef. : \textbf{7}\end{minipage}
\par\smallskip{\color{onbo}\rule{\linewidth}{1pt}}

\subsection*{Partie A — Évaluation des ressources \hfill (15 points)}

\noindent\textbf{Exercice 1.} \hfill\textit{(3 points)}\\
Un grenier contient $5$ sacs de maïs jaune et $3$ sacs de maïs blanc, indiscernables au
toucher. Un commerçant tire au hasard et simultanément $3$ sacs. Soit $X$ la variable
aléatoire égale au nombre de sacs de maïs blanc obtenus.
\begin{enumerate}[label=\arabic*)]
\item Déterminer la loi de probabilité de $X$. \hfill\textit{(1,5 pt)}
\item Calculer l'espérance mathématique $E(X)$. \hfill\textit{(0,75 pt)}
\item Calculer la probabilité d'obtenir au plus un sac de maïs blanc. \hfill\textit{(0,75 pt)}
\end{enumerate}

\smallskip
\noindent\textbf{Exercice 2.} \hfill\textit{(3 points)}\\
\begin{enumerate}[label=\arabic*)]
\item Résoudre dans $\Z^2$ l'équation $11x-8y=3$. \hfill\textit{(1,5 pt)}
\item Déterminer le reste de la division euclidienne de $11^{2025}$ par $7$. \hfill\textit{(1,5 pt)}
\end{enumerate}

\smallskip
\noindent\textbf{Exercice 3.} \hfill\textit{(4 points)}\\
Soit $f$ la fonction définie sur $]0,+\infty[$ par $f(x)=x-2+\dfrac{\ln x}{x}$, de courbe $(\mathcal{C})$.
\begin{enumerate}[label=\arabic*)]
\item Déterminer les limites de $f$ en $0^+$ et en $+\infty$. Montrer que la droite
$(D):y=x-2$ est asymptote à $(\mathcal{C})$ en $+\infty$. \hfill\textit{(1,25 pt)}
\item Montrer que $f'(x)=\dfrac{x^2+1-\ln x}{x^2}$ et en déduire le sens de variation de $f$. \hfill\textit{(1,5 pt)}
\item À l'aide d'une intégration par parties, calculer $\displaystyle I=\int_1^{\mathrm{e}}\frac{\ln x}{x}\,\dd x$,
puis en déduire $\displaystyle\int_1^{\mathrm{e}}f(x)\,\dd x$. \hfill\textit{(1,25 pt)}
\end{enumerate}

\smallskip
\noindent\textbf{Exercice 4.} \hfill\textit{(5 points)}\\
Le plan complexe est muni d'un repère orthonormé direct $(O\,;\,\vec u,\vec v)$.
On considère les points $A$, $B$, $C$ d'affixes respectives
$z_A=2$, $z_B=3+i\sqrt3$ et $z_C=1+i\sqrt3$.
\begin{enumerate}[label=\arabic*)]
\item Calculer le module et un argument de $z_B-z_A$ et de $z_C-z_A$. \hfill\textit{(1,5 pt)}
\item En déduire la nature du triangle $ABC$. \hfill\textit{(1 pt)}
\item Déterminer l'affixe du point $G$, barycentre des points pondérés
$(A,2)$, $(B,1)$, $(C,1)$. \hfill\textit{(1 pt)}
\item Soit $S$ la similitude directe de centre $A$, de rapport $2$ et d'angle $\frac\pi3$.
Donner l'écriture complexe de $S$ et déterminer l'affixe de l'image $B'$ de $B$ par $S$. \hfill\textit{(1,5 pt)}
\end{enumerate}

\subsection*{Partie B — Évaluation des compétences \hfill (5 points)}

\noindent\textit{Monsieur Ngono exploite un champ de maïs à Bafia. Il sollicite ton aide, en
tant qu'élève de Terminale~C, pour préparer sa campagne. Il dispose de \SI{240}{\metre} de
clôture et veut délimiter une parcelle rectangulaire \emph{adossée à une rivière} (le côté
longeant la rivière ne se clôture pas). Sur cette parcelle, son rendement est de
\SI{3,5}{\tonne} la première année puis augmente de \SI{6}{\percent} chaque année grâce à
l'amélioration des sols. Enfin, à la récolte, $4\,\%$ des épis sont attaqués par un
charançon ; un contrôleur prélève $5$ épis au hasard dans un grand stock.}

\begin{enumerate}[label=\textbf{Tâche \arabic*.},leftmargin=2.4em]
\item Déterminer les dimensions de la parcelle rectangulaire d'\emph{aire maximale} qu'il
peut clôturer, ainsi que cette aire. \hfill\textit{(1,5 pt)}
\item Déterminer l'année à partir de laquelle son rendement annuel dépassera \SI{5}{\tonne}. \hfill\textit{(1,5 pt)}
\item Déterminer la probabilité qu'\emph{au moins un} des $5$ épis prélevés soit attaqué. \hfill\textit{(1,5 pt)}
\end{enumerate}
\par\smallskip{\footnotesize\itshape Présentation générale : 0,5 point.}

% ════════════════════════════════════════════════════
\corrige{du sujet 10}

\noindent\textbf{Partie A — Exercice 1.}
1) $\binom{8}{3}=56$ tirages équiprobables. $X\in\{0,1,2,3\}$ et
$P(X{=}k)=\dfrac{\binom3k\binom5{3-k}}{56}$ :
$P(0)=\frac{10}{56}=\frac5{28},\ P(1)=\frac{30}{56}=\frac{15}{28},\ P(2)=\frac{15}{56},\ P(3)=\frac1{56}$.
(Vérification : $\frac{10+30+15+1}{56}=1$.)
2) $E(X)=\dfrac{0\cdot10+1\cdot30+2\cdot15+3\cdot1}{56}=\dfrac{63}{56}=\dfrac98=1{,}125$.
3) $P(X\le1)=\frac{10}{56}+\frac{30}{56}=\frac{40}{56}=\frac57$.

\noindent\textbf{Exercice 2.}
1) $11\wedge8=1$ donc l'équation a des solutions. $11\times3-8\times4=33-32=1$, donc
$11\times9-8\times12=3$ : une solution particulière est $(9\,;\,12)$. Par différence,
$11(x-9)=8(y-12)$ ; comme $8\wedge11=1$, le théorème de Gauss donne $8\mid x-9$, d'où
$x=9+8k$ et $y=12+11k$, $k\in\Z$.
2) Modulo $7$ : $11\equiv4$. Les puissances de $4$ : $4^1\equiv4$, $4^2=16\equiv2$,
$4^3\equiv8\equiv1\ [7]$, de période $3$. Comme $2025=3\times675$,
$11^{2025}\equiv4^{2025}\equiv(4^3)^{675}\equiv1\ [7]$. Le reste est $\mathbf{1}$.

\noindent\textbf{Exercice 3.}
1) En $0^+$ : $x-2\to-2$ et $\frac{\ln x}{x}\to-\infty$, donc $f(x)\to-\infty$.
En $+\infty$ : $\frac{\ln x}{x}\to0$, donc $f(x)\to+\infty$. De plus
$f(x)-(x-2)=\frac{\ln x}{x}\to0$ en $+\infty$ : la droite $(D):y=x-2$ est asymptote à $(\mathcal{C})$.
2) $\Big(\frac{\ln x}{x}\Big)'=\frac{\frac1x\cdot x-\ln x}{x^2}=\frac{1-\ln x}{x^2}$, donc
$f'(x)=1+\frac{1-\ln x}{x^2}=\frac{x^2+1-\ln x}{x^2}$. Pour $x>0$, on a $\ln x\le x-1<x^2+1$,
donc le numérateur $x^2+1-\ln x>0$ : $f'(x)>0$ et $f$ est \emph{strictement croissante} sur $]0,+\infty[$.
3) IPP avec $u=\ln x$, $\dd u=\frac{\dd x}{x}$ : $I=\int_1^{\mathrm e}\ln x\cdot\frac1x\dd x=\big[\tfrac12(\ln x)^2\big]_1^{\mathrm e}=\tfrac12$.
(On peut aussi poser $t=\ln x$.) Alors
$\int_1^{\mathrm e}f(x)\dd x=\int_1^{\mathrm e}(x-2)\dd x+I=\big[\tfrac{x^2}2-2x\big]_1^{\mathrm e}+\tfrac12
=\big(\tfrac{\mathrm e^2}2-2\mathrm e\big)-\big(\tfrac12-2\big)+\tfrac12=\dfrac{\mathrm e^2}2-2\mathrm e+2$.

\noindent\textbf{Exercice 4.}
1) $z_B-z_A=1+i\sqrt3=2\big(\tfrac12+i\tfrac{\sqrt3}2\big)=2\mathrm e^{i\pi/3}$ : module $2$, argument $\frac\pi3$.
$z_C-z_A=-1+i\sqrt3=2\big(-\tfrac12+i\tfrac{\sqrt3}2\big)=2\mathrm e^{i2\pi/3}$ : module $2$, argument $\frac{2\pi}3$.
2) $AB=|z_B-z_A|=2$ et $AC=|z_C-z_A|=2$, donc $AB=AC$ : triangle isocèle en $A$. De plus
$\arg\frac{z_C-z_A}{z_B-z_A}=\frac{2\pi}3-\frac\pi3=\frac\pi3$, donc $\widehat{BAC}=\frac\pi3$ :
le triangle $ABC$ est \emph{équilatéral}.
3) Masse totale $2+1+1=4$. $z_G=\dfrac{2z_A+z_B+z_C}{4}=\dfrac{4+(3+i\sqrt3)+(1+i\sqrt3)}{4}=\dfrac{8+2i\sqrt3}{4}=2+\tfrac{\sqrt3}2 i$.
4) $S(z)=a(z-z_A)+z_A$ avec $a=2\mathrm e^{i\pi/3}=2\big(\tfrac12+i\tfrac{\sqrt3}2\big)=1+i\sqrt3$.
Donc $S(z)=(1+i\sqrt3)(z-2)+2$. Image de $B$ :
$z_{B'}=(1+i\sqrt3)(z_B-2)+2=(1+i\sqrt3)(1+i\sqrt3)+2=(1+i\sqrt3)^2+2$.
Or $(1+i\sqrt3)^2=1+2i\sqrt3-3=-2+2i\sqrt3$, d'où $z_{B'}=2i\sqrt3$.

\smallskip
\noindent\textbf{Partie B — Situation.}
\emph{Tâche 1.} Notons $x$ la mesure (en m) de chacun des deux côtés perpendiculaires à la
rivière et $y$ le côté parallèle à la rivière. La clôture vérifie $2x+y=240$, soit
$y=240-2x$ (avec $0<x<120$). L'aire vaut $A(x)=x(240-2x)=240x-2x^2$, et
$A'(x)=240-4x=0$ donne $x=60$. Comme $A''(x)=-4<0$, c'est un maximum. Dimensions :
$x=\SI{60}{\metre}$ et $y=\SI{120}{\metre}$ ; aire maximale $A=60\times120=\SI{7200}{\metre\squared}$.

\emph{Tâche 2.} Le rendement de l'année $n$ est $u_n=3{,}5\times1{,}06^{\,n-1}$ (suite
géométrique de raison $1{,}06$). On résout $u_n>5$, soit $1{,}06^{\,n-1}>\frac{5}{3{,}5}=\frac{10}{7}$,
d'où $n-1>\dfrac{\ln(10/7)}{\ln1{,}06}\approx\dfrac{0{,}3567}{0{,}0583}\approx6{,}12$, soit $n>7{,}12$.
Le rendement dépasse \SI{5}{\tonne} à partir de la \textbf{8\textsuperscript{e} année}
($u_8=3{,}5\times1{,}06^{7}\approx\SI{5,26}{\tonne}$, alors que $u_7\approx\SI{4,96}{\tonne}$).

\emph{Tâche 3.} Les $5$ prélèvements sont assimilés à des tirages indépendants (grand stock),
chacun « attaqué » avec probabilité $0{,}04$. Donc
$P(\text{au moins un attaqué})=1-(0{,}96)^5\approx1-0{,}8154=\mathbf{0{,}1846}$, soit environ \SI{18,5}{\percent}.
